\section{Proof of the global atlas bounds}\label{app:global-atlas}

\subsection*{Transition chart}
For $A>0$ and $r\in\mathbb N_0$, elementary differentiation shows
\[
 \sup_{0<t\le t_0}e^{-A/t}t^{-r}=\mathcal H_r(A,t_0),
\]
with \eqref{eq:Hmax} for $r=0$ and with the unique unrestricted maximum
at $t=A/r$ when $r\ge1$.  Applying this identity separately to the two
terms on the right of \eqref{eq:Fa-bound} gives
\[
 \sup_{0<t\le t_0}\sup_{0\le\theta\le R}|q_t^{(m)}(\theta)|
 \le \widehat Q_m(t_0,R).
\]
For large $k$, $A_k(R)/r>t_0$, so the summands reduce to a polynomial in
$a_k$ times $e^{-A_k(R)/t_0}$.  Successive terms therefore admit the same
kind of geometric tail as Appendix~\ref{app:image-tails}.  If
$\widehat Q_0<1$, Fa\`a di Bruno's formula and the Euler-product estimate
give \eqref{eq:global-transition}.

\subsection*{Cut chart}
Write
\begin{equation}\label{eq:Jbt}
 J_{b,t}(\phi)=\frac{b^2}{t}\int_0^1
 \cosh\frac{sb\phi}{t}\,\dd s-\cosh\frac{b\phi}{t}.
\end{equation}
Since
\[
 b\sinh\frac{b\phi}{t}-\phi\cosh\frac{b\phi}{t}
 =\phi J_{b,t}(\phi),
\]
the apparent singularity at $\phi=0$ is removed before differentiation.
The heat kernel on the cut chart is exactly
\begin{equation}\label{eq:cut-sum-regular}
 p_t(\pi-\phi)=2C_t\sum_{j\ge0}e^{-b_j^2/(2t)}
 e^{-\phi^2/(2t)}\frac{\phi}{\sin\phi}J_{b_j,t}(\phi).
\end{equation}

We now give fully explicit derivative bounds for its three factors.  The
Hermite-polynomial formula implies
\begin{equation}\label{eq:Ecut}
 \left|\partial_\phi^r e^{-\phi^2/(2t)}\right|
 \le \mathcal E_r(t,\delta):=
 r!\sum_{q=0}^{\lfloor r/2\rfloor}
 \frac{\delta^{r-2q}}{2^q q!(r-2q)!}\,t^{-r+q}.
\end{equation}
Let $c_n$ be the positive Taylor coefficients of $x/\sin x$,
\[
 \frac{x}{\sin x}=\sum_{n\ge0}c_nx^{2n},\qquad
 c_0=1,\qquad
 c_n=\frac{(-1)^n2(1-2^{2n-1})B_{2n}}{(2n)!}\quad(n\ge1).
\]
Here $B_{2n}$ denotes the Bernoulli numbers.
For $0<\delta<\pi$, termwise differentiation gives the convergent bound
\begin{equation}\label{eq:Scut}
 \left|\partial_\phi^s\frac{\phi}{\sin\phi}\right|
 \le \mathcal S_s(\delta):=
 \sum_{n\ge\lceil s/2\rceil}c_n
 \frac{(2n)!}{(2n-s)!}\delta^{2n-s}.
\end{equation}
Finally, differentiation under the integral in \eqref{eq:Jbt} yields
\begin{equation}\label{eq:Jcut}
 |\partial_\phi^vJ_{b,t}(\phi)|
 \le \mathcal J_v(b;t,\delta):=
 e^{b\delta/t}\left(\frac bt\right)^v
 \left(1+\frac{b^2}{t(v+1)}\right).
\end{equation}

Define
\begin{equation}\label{eq:Dcut}
 D_m^{\rm cut}(t,\delta)=2C_t\sum_{j\ge0}e^{-b_j^2/(2t)}
 \sum_{r+s+v=m}\frac{m!}{r!s!v!}
 \mathcal E_r(t,\delta)\mathcal S_s(\delta)
 \mathcal J_v(b_j;t,\delta).
\end{equation}
The Gaussian $e^{-b_j^2/(2t)+b_j\delta/t}$ makes the series and all its
derivatives summable.  Equations~\eqref{eq:Ecut}--\eqref{eq:Jcut} and the
Leibniz rule prove
\begin{equation}\label{eq:pcut-derivative}
 \sup_{0\le\phi\le\delta}
 |\partial_\phi^m p_t(\pi-\phi)|\le D_m^{\rm cut}(t,\delta).
\end{equation}

For the lower bound, the inequality used in Appendix~\ref{app:cut} gives
\[
 J_{b,t}(\phi)
 =\cosh\frac{b\phi}{t}
 \left(\frac{b}{\phi}\tanh\frac{b\phi}{t}-1\right)
 \ge \frac{b^2}{t+b\delta}-1.
\]
Under $\pi^2>t+\pi\delta$, every summand in
\eqref{eq:cut-sum-regular} is positive.  Retaining the principal pair and
using $\phi/\sin\phi\ge1$ gives
\begin{equation}\label{eq:Lcut}
 p_t(\pi-\phi)\ge L_{\rm cut}(t,\delta):=
 2C_t e^{-(\pi^2+\delta^2)/(2t)}
 \left(\frac{\pi^2}{t+\pi\delta}-1\right)>0.
\end{equation}
Set
\begin{equation}\label{eq:Ccut}
 C_m^{\rm cut}(t,\delta)=
 \sum_{\ell=1}^m\frac{(\ell-1)!}{L_{\rm cut}(t,\delta)^\ell}
 \Bell_{m,\ell}\left(D_1^{\rm cut},\ldots,
 D_{m-\ell+1}^{\rm cut}\right).
\end{equation}
Fa\`a di Bruno's formula applied to $\log p_t$ proves
\eqref{eq:global-cut}.  Since $\partial_\theta=-\partial_\phi$, the absolute
derivative bounds are unchanged.

\subsection*{Antipodal asymptotics}
Taking the removable limit in \eqref{eq:cut-sum-regular} gives the exact sum
\begin{equation}\label{eq:antipode-exact}
 p_t(\pi)=2C_t\sum_{j\ge0}e^{-b_j^2/(2t)}
 \left(\frac{b_j^2}{t}-1\right).
\end{equation}
The $j=0$ term and the gap $b_1^2-b_0^2=8\pi^2$ prove
\eqref{eq:antipode-asymptotic}.  Expanding the principal summand of
\eqref{eq:cut-sum-regular} through order $\phi^2$, then differentiating
$-\log p_t$, gives
\[
 s_t''(\pi)=-\frac{\pi^2}{3t^2}+O(t^{-1});
\]
the later pairs are smaller by $O(e^{-4\pi^2/t})$.  This proves
\eqref{eq:antipode-hessian} and completes Theorem~\ref{thm:global-atlas}.
